%---------------------------------------------------------------------
% APPENDIX --- FORMULA SHEET AND EXAM METHOD
%---------------------------------------------------------------------
\chapter{Formula Sheet}

Everything the nine units require, on four pages. Learn the \emph{shape} of
each rule and where it applies; the arithmetic is the easy part.

%---------------------------------------------------------------------
\section{Probability (Units 1--2)}

\begin{center}
\small
\begin{tabular}{@{}L{5.2cm} L{4.9cm} L{5.1cm}@{}}
\toprule
\textbf{Rule} & \textbf{Formula} & \textbf{Use it when} \\
\midrule
Complement & $\prob{A'}=1-\prob{A}$ & ``at least one'', ``not'' \\
Addition (general) & $\prob{A\cup B}=\prob{A}+\prob{B}-\prob{A\cap B}$ & ``or'', events overlap \\
Addition (exclusive) & $\prob{A\cup B}=\prob{A}+\prob{B}$ & ``or'', $A\cap B=\varnothing$ \\
Multiplication (indep.) & $\prob{A\cap B}=\prob{A}\prob{B}$ & ``and'', with replacement \\
Multiplication (dep.) & $\prob{A\cap B}=\prob{A}\prob{B\mid A}$ & ``and'', without replacement \\
Conditional & $\prob{A\mid B}=\dfrac{\prob{A\cap B}}{\prob{B}}$ & ``given that'' \\
Independence test & $\prob{A\cap B}=\prob{A}\prob{B}$ & to prove or disprove it \\
\bottomrule
\end{tabular}
\end{center}

\medskip
\begin{center}
\small
\begin{tabular}{@{}L{4.0cm} L{6.3cm} L{4.9cm}@{}}
\toprule
\textbf{Distribution} & \textbf{Formula} & \textbf{Recognise it by} \\
\midrule
Binomial &
$P(X=x)=\dbinom{n}{x}p^{x}(1-p)^{n-x}$, $E(X)=np$ &
$n$ fixed trials, constant $p$, \emph{with} replacement \\[6pt]
Hypergeometric &
$P(X=x)=\dfrac{\binom{K}{x}\binom{N-K}{n-x}}{\binom{N}{n}}$ &
sampling from a finite bag \emph{without} replacement \\[6pt]
Any discrete &
$\sum_x P(X=x)=1$, \quad $E(X)=\sum_x x\,P(X=x)$ &
always --- use as a check \\
\bottomrule
\end{tabular}
\end{center}

%---------------------------------------------------------------------
\section{Equations, inequalities, straight lines (Units 3--4)}

\begin{center}
\small
\begin{tabular}{@{}L{5.0cm} L{10.6cm}@{}}
\toprule
Quadratic formula & $x=\dfrac{-b\pm\sqrt{b^{2}-4ac}}{2a}$; discriminant
$b^{2}-4ac>0$ two roots, $=0$ one, $<0$ none real \\
Completing the square & $x^{2}+bx=\left(x+\tfrac b2\right)^{2}-\tfrac{b^{2}}{4}$ \\
Inequalities & multiplying or dividing by a \emph{negative} reverses the sign \\
Absolute value (equation) & $|u|=|v| \iff u=v$ or $u=-v$ \\
Absolute value ($\le$) & $|u|\le a \iff -a\le u\le a$ \quad (one interval) \\
Absolute value ($\ge$) & $|u|\ge a \iff u\ge a$ or $u\le -a$ \quad (two intervals) \\
Distance & $d=\sqrt{(x_2-x_1)^{2}+(y_2-y_1)^{2}}$ \\
Midpoint & $M=\left(\dfrac{x_1+x_2}{2},\;\dfrac{y_1+y_2}{2}\right)$ \\
Slope & $m=\dfrac{y_2-y_1}{x_2-x_1}$ \\
Point--slope / slope--intercept & $y-y_1=m(x-x_1)$; \quad $y=mx+c$ \\
Parallel / perpendicular & $m_1=m_2$; \quad $m_1m_2=-1$ (not for vertical lines) \\
Linear cost & $C(x)=(\text{variable per unit})x+(\text{fixed cost})$ \\
Equilibrium & set demand $=$ supply and solve for $q$, then substitute for $p$ \\
\bottomrule
\end{tabular}
\end{center}

%---------------------------------------------------------------------
\section{Matrices and determinants (Units 5--6)}

\begin{center}
\small
\begin{tabular}{@{}L{5.0cm} L{10.6cm}@{}}
\toprule
Conformability & $A_{m\times n}B_{n\times p}$ is defined and is $m\times p$ \\
Transpose rules & $(A^{t})^{t}=A$, \ $(A+B)^{t}=A^{t}+B^{t}$, \
                  $\boxed{(AB)^{t}=B^{t}A^{t}}$ \\
Inverse rules & $(A^{-1})^{-1}=A$, \ $\boxed{(AB)^{-1}=B^{-1}A^{-1}}$ \\
$2\times2$ determinant & $\dt{a&b\\c&d}=ad-bc$ \\
$2\times2$ inverse & $\dfrac{1}{ad-bc}\mat{d&-b\\-c&a}$ \\
Cofactor & $C_{ij}=(-1)^{i+j}M_{ij}$; sign pattern $\mat{+&-&+\\-&+&-\\+&-&+}$ \\
Determinant & $\det A=\sum_j a_{ij}C_{ij}$ along any row or column \\
Adjoint & $\operatorname{adj}A=\left[C_{ij}\right]^{t}$ --- \emph{transpose} the
          cofactor matrix \\
Inverse & $A^{-1}=\dfrac{1}{\det A}\operatorname{adj}A$, requires $\det A\neq0$ \\
Cramer's rule & $x_k=\dfrac{D_k}{D}$, $D=\det A$, $D_k$ replaces column $k$
                 by $b$ \\
$\det=0$ means & no unique solution: either none, or infinitely many \\
Product rule for $\det$ & $\det(AB)=\det A\cdot\det B$ \\
\bottomrule
\end{tabular}
\end{center}

\begin{keyidea}[Determinant shortcuts]
Expand along the row or column with most zeros. A common factor in a row or
column comes outside. Two identical or proportional rows (or columns) force
$\det=0$. Adding a multiple of one row to another leaves $\det$ unchanged.
\end{keyidea}

%---------------------------------------------------------------------
\section{Calculus (Units 7--9)}

\begin{center}
\small
\begin{tabular}{@{}L{4.4cm} L{5.4cm} L{5.4cm}@{}}
\toprule
\textbf{Rule} & \textbf{Formula} & \textbf{Note} \\
\midrule
Power    & $\dd{}{x}x^{n}=nx^{n-1}$ & rewrite roots and reciprocals as powers first \\
Constant multiple & $\dd{}{x}\left[cf\right]=cf'$ & \\
Sum      & $(u+v)'=u'+v'$ & \\
Product  & $(uv)'=u'v+uv'$ & \\
Quotient & $\left(\dfrac uv\right)'=\dfrac{u'v-uv'}{v^{2}}$ & derivative of the \emph{top} first \\
Chain    & $\dd{}{x}f(g(x))=f'(g(x))\,g'(x)$ & outer first, then $\times$ inner \\
Exponential & $\dd{}{x}e^{u}=e^{u}\,u'$ & \\
Logarithm & $\dd{}{x}\ln u=\dfrac{u'}{u}$ & \\
\bottomrule
\end{tabular}
\end{center}

\medskip
\begin{center}
\small
\begin{tabular}{@{}L{5.0cm} L{10.6cm}@{}}
\toprule
Limits at infinity & divide by the highest power of $x$ in the denominator \\
$\tfrac00$ form & factorise and cancel, then substitute again \\
Limit exists at $a$ & $\displaystyle\lim_{x\to a^{-}}f=\lim_{x\to a^{+}}f$ \\
Partial derivative & differentiate in one variable, treat the other as a constant \\
Clairaut & $f_{xy}=f_{yx}$ --- use as a free check on mixed partials \\
Stationary point & $f'(x)=0$ \ (or $f_x=f_y=0$ in two variables) \\
Second-derivative test (1 var) & $f''<0$ max, $f''>0$ min, $f''=0$ inconclusive \\
Second-derivative test (2 var) & $D=f_{xx}f_{yy}-(f_{xy})^{2}$; $D>0$ with
  $f_{xx}<0$ max, with $f_{xx}>0$ min; $D<0$ saddle; $D=0$ inconclusive \\
Inflection point & $f''=0$ \emph{and} $f''$ changes sign \\
Revenue & $R(x)=x\cdot p(x)$ \\
Profit & $P(x)=R(x)-C(x)$; maximised where $MR=MC$ and $P''<0$ \\
Average cost & $\bar{C}(x)=\dfrac{C(x)}{x}$; at its minimum, $MC=\bar{C}$ \\
\bottomrule
\end{tabular}
\end{center}

%---------------------------------------------------------------------
\chapter{Method in the Exam Hall}

\section{The eight recurring mistakes}

Collected from every \emph{Where marks are lost} box in this document, in the
order the units come:

\begin{enumerate}
  \item \textbf{Using $\prob{A}+\prob{B}$ for overlapping events.} The simple
        addition rule needs mutual exclusivity. Otherwise subtract the
        intersection.
  \item \textbf{Forgetting the count in ``exactly one''.} $\binom{n}{1}$ counts
        \emph{which} one. ``All'' and ``none'' need no such factor.
  \item \textbf{Binomial where the problem says ``without replacement''.} That
        phrase means hypergeometric, and it changes the answer materially.
  \item \textbf{Not reversing the inequality} when multiplying or dividing by a
        negative number.
  \item \textbf{$(AB)^{t}=A^{t}B^{t}$ and $(AB)^{-1}=A^{-1}B^{-1}$.} Both
        reverse the order. Both are worth a whole question.
  \item \textbf{Forgetting to transpose the cofactor matrix} when forming the
        adjoint. Check an \emph{off-diagonal} entry of $AA^{-1}$, not just a
        diagonal one.
  \item \textbf{Not computing $D$ before applying Cramer's rule.} If $D=0$ the
        rule does not apply and the correct answer is a statement about the
        system.
  \item \textbf{Stopping at $f'=0$.} The second-derivative test is separately
        marked, and it is what distinguishes maximum from minimum from saddle.
\end{enumerate}

\section{Six habits worth the seconds they cost}

\begin{description}
  \item[Substitute your answer back.] Nearly every solution in this document
        ends with a check line. It costs fifteen seconds and catches sign
        errors, which are the most common and the most expensive.
  \item[Tabulate before you compute.] The probability questions in Unit 1 and
        the matrix word problems in Unit 5 become trivial once the data is in a
        two-way table, and nearly impossible without one.
  \item[Define your variable in writing.] ``Let $n$ = the number of \$0.60
        decreases.'' Most marks lost in Unit 9 word problems are lost before
        any calculus starts.
  \item[Simplify algebraically first.] $A+2(B-A)$ is $2B-A$; a common factor in
        the numerator of a partial derivative often collapses the whole
        expression. Do the algebra before the arithmetic.
  \item[Sanity-check the units and the sign.] A negative quantity, a
        probability above 1, an average cost below marginal cost --- each is a
        signal to go back, and saying so in your script earns credit even if
        you cannot find the slip.
  \item[Show the method even when stuck.] AIOU awards method marks. A correct
        formula with a wrong number scores; a right number with no working
        often does not.
\end{description}

\section{If revision time is short}

Units 4, 5, 6, 7 and 9 carry most of the marks and all of the long questions.
Within those, the five highest-yield techniques are:

\begin{enumerate}
  \item Cramer's rule on a $3\times3$ system, including the $D=0$ case.
  \item Inverse of a $3\times3$ matrix by cofactors, with the adjoint transposed.
  \item Product, quotient and chain rules --- especially chain inside quotient.
  \item Optimising profit, revenue or average cost, with the second-order test.
  \item Finding a line's equation from two points, or from a point and a
        perpendicularity condition.
\end{enumerate}

Work each of these from the relevant unit in this document with the solution
covered. Five techniques, done cold and correctly, are worth more than
nine units read passively.
